\documentclass[11pt]{article}

\usepackage[margin=1in]{geometry}
\usepackage{amsmath,amssymb,mathtools}
\usepackage{booktabs}
\usepackage{microtype}
\usepackage{xcolor}
\usepackage{hyperref}
\usepackage{url}
\makeatletter
\g@addto@macro\UrlBreaks{\do\a\do\b\do\c\do\d\do\e\do\f\do\g\do\h\do\i\do\j\do\k\do\l\do\m\do\n\do\o\do\p\do\q\do\r\do\s\do\t\do\u\do\v\do\w\do\x\do\y\do\z\do\A\do\B\do\C\do\D\do\E\do\F\do\G\do\H\do\I\do\J\do\K\do\L\do\M\do\N\do\O\do\P\do\Q\do\R\do\S\do\T\do\U\do\V\do\W\do\X\do\Y\do\Z\do\0\do\1\do\2\do\3\do\4\do\5\do\6\do\7\do\8\do\9}
\makeatother
\usepackage{longtable}

\hypersetup{
    colorlinks=true,
    linkcolor=blue,
    citecolor=blue,
    urlcolor=blue,
    pdftitle={Computing Machinery and Understanding},
    pdfauthor={M. Drake Stapleton}
}


\title{
    \textbf{Computing Machinery and Understanding}\\[0.5em]
    \large The Turing as a Measure Emerging from AIEN
}
\author{M. Drake Stapleton}
\date{October 2, 2026\\v1.7 preprint (not peer reviewed)}

\begin{document}

\maketitle

\noindent\textit{Companion to the formal preprint ``Computing Machinery and Understanding: From the Imitation Game to the Turing'' (v1.13), which gives the definitions, proofs, measurement protocol, and falsification criteria. This paper gives the engineering history the measure emerged from.}

\begin{abstract}
I did not begin this work by trying to invent a new unit of intelligence.

I began by trying to build a machine whose internal claims could be separated from evidence, whose ideas could be separated from their physical implementations, whose failures could remain part of its history, and whose improvements could be tested against reality before they became authoritative.

That work became AIEN, Omega, FORGE, AEGIS, ARGUS, Cortex, and the resident reaction architecture now implemented across the AIEN repositories.

While working inside that layer of abstraction, I encountered a measurement problem.

The system could generate alternatives. It could verify them. It could measure them on real hardware. It could preserve failed hypotheses. It could discover recurring abstractions. It could promote a verified improvement into a new generation. It could survive injected failures and reconstruct itself from durable state.

But I still lacked a satisfactory answer to a simpler question:

\begin{quote}
\textbf{How much did the machine actually learn?}
\end{quote}

Tokens per second did not answer it. Accuracy did not answer it. FLOPs did not answer it. Power consumption did not answer it. Even successful self-improvement did not answer it, because an improvement can be narrow, expensive, memorized, or dependent on complexity that exceeds what it explains.

The quantity I arrived at is a direct consequence of the architecture I had already built.

For a declared baseline \(B\), candidate explanation \(M\), sealed held-out observations \(D\), and measurement profile \(P\), I define the Turing gain

\[
\mathcal{T}(M;B,D,P)
=
DL(B,D)-DL(M,D),
\]

where

\[
DL(M,D)=L(M)+L(D\mid M).
\]

One \textbf{Turing}, \(1T\), is one bit of net held-out explanatory compression under the declared protocol.

The important claim of this paper is therefore not that I have renamed the bit.

It is that, while building an architecture in which explanations, measurements, realizations, authority, provenance, and physical cost were already separate objects, I found that \textbf{net held-out description-length reduction is the quantity that connects them}.

AIEN already implements much of the machinery required to test that proposition. This paper documents that implementation, the evidence already produced, the places where the architecture failed and corrected itself, and the first profile-bound idealized Turing result, +2,559,679.825~\(T_{\mathrm{ideal}}\) under profile V0, recorded 2026-09-29: a candidate explanation of Omega search traces scored +2,559,679.825~\(T_{\mathrm{ideal}}\) under profile V0 against its baseline on sealed held-out seeds, about 0.44~\(T_{\mathrm{ideal}}\) per held-out observation of specific Turing gain. That result is a theoretical description length, not yet a physically encoded file, and no energy was measured in the run, so no Turing yield is claimed. What the unit does not measure, I state plainly: semantic grounding, causal direction, agency, values, alignment, and consciousness are outside its scope. The Turing is a thermometer for one quantity, net held-out explanatory compression, and the claim is only that this quantity is worth measuring well.

\medskip

\noindent\emph{Turing taught us to replace an unanswerable question with an operational one. My question is whether a machine can discover explanations that make unseen reality less surprising, and whether I can measure the physical cost of that discovery.}
\end{abstract}

\section{I Did Not Start With the Turing}

My path into this problem ran in the opposite direction from the usual benchmark paper.

I was not trying to create another score for artificial intelligence.

I was trying to create a computational system in which I could tell, after the fact, \emph{why something happened, what evidence caused it, which physical realization produced it, whether that realization had authority to act, whether the result survived verification, and whether an improvement could be reproduced later}.

The source code is public:

\begin{itemize}
    \item AIEN/Omega: \url{https://github.com/aien-dev/omega}
    \item AIEN architecture: \url{https://github.com/aien-dev/aien-architecture}
    \item FORGE: \url{https://github.com/aien-dev/physics} (repository retains the historical PHYSICS name)
    \item AIENOS: \url{https://github.com/aien-dev/aienos}
    \item Sovereign core: \url{https://github.com/aien-dev/aien-sovereign-core}
    \item Benchmarks: \url{https://github.com/aien-dev/benchmarks}
\end{itemize}

The abstraction I eventually converged on separates four things that are usually collapsed together:

\[
\text{meaning}
\neq
\text{realization}
\neq
\text{authority}
\neq
\text{evidence}.
\]

Omega defines semantics and synthesizes candidates.

FORGE realizes and measures physical realizations.

AEGIS certifies generated code, including crumbs, according to policy and verification evidence. It changes nothing, decides no policy, and is not part of the generation system.

ARGUS observes defensive events without becoming an authority itself.

The World holds operational state.

Cortex retains epistemic memory.

The evidence system records what actually occurred.

That separation turned out to matter more than I expected.

Once a semantic operation is no longer identical to one implementation, the system can ask:

\begin{quote}
What other realization preserves the same meaning?
\end{quote}

Once a hypothesis is no longer identical to authority, it can ask:

\begin{quote}
What could I try without allowing the hypothesis to approve itself?
\end{quote}

Once an outcome is no longer identical to a claim, it can ask:

\begin{quote}
What evidence actually supports this?
\end{quote}

And once a result is no longer identical to success, a further question appears:

\begin{quote}
Was the new explanation actually simpler and more predictive than what I had before?
\end{quote}

That final question is where the Turing came from.

\section{The First Piece Already Existed: Omega Was Discovering Compression}

Before I formulated the Turing, Omega already contained an implemented abstraction-discovery system.

The implementation is public:

\begin{itemize}
    \item \url{https://github.com/aien-dev/omega/blob/4f55f45ed13f63cf8171bebf223c6258f8351ab9/src/omega_discovery.c}
    \item \url{https://github.com/aien-dev/omega/blob/1ee43fbf1c7d588734b80bd0743e5d054aea750b/src/omega_discovery.h}
\end{itemize}

The experiment is deliberately small.

Omega constructs a corpus of independently verified programs containing a repeated structure. In the canonical qualification corpus, that repeated structure corresponds to the program

\[
2x+1.
\]

The implementation scans realized programs for recurring instruction slices. A candidate must occur in more than one program and must be nontrivial.

The important part is the scoring rule. The source computes a net structural compression score conceptually equivalent to

\[
\text{compression score}
=
\text{occurrences}\times(L-1)-(L+1),
\]

where replacing a recurring sequence saves instructions at each occurrence while introducing the abstraction itself also has a cost.

Thus the abstraction is useful only when

\[
\text{structure removed}
>
\text{cost of explanation}.
\]

This was already a primitive Minimum Description Length calculation in the sense of Rissanen and Grünwald~\cite{rissanen1978,grunwald2007}.

The qualification receipt is public:

\begin{quote}
\url{https://github.com/aien-dev/omega/blob/8e0a50a45814c5e3b427efd36a83a1e2e79a7500/evidence/omega_library_discovery_qualification_receipt.json}
\end{quote}

That receipt records:

\begin{itemize}
    \item autonomous recurring-subsequence mining;
    \item \(+7\) instructions of net structural compression;
    \item semantic preservation after refactoring;
    \item admission checks through Omega's V0 to V2 ladder;
    \item reduction in held-out synthesis exploration from \(54\) candidates to \(21\);
    \item all ten M11 qualification gates passing.
\end{itemize}

This history matters.

I did not begin with

\[
L(M)+L(D\mid M)
\]

and decide to force Omega into it.

I was already building machinery that admitted a reusable abstraction only when the abstraction paid for itself.

The Turing is the generalization of that principle from program structure to predictive explanation.

\section{From Static Abstraction to Continuation}

Structural compression alone is not enough.

A machine can discover a useful function without being a system that continues to investigate, test, and alter itself.

The next architectural problem I worked on was therefore causal continuation.

The architecture stopped treating cognition as a central loop that explicitly called every subsystem in order.

Instead, state changes made dependent reactions ready.

The relevant runtime work lives in:

\begin{quote}
\url{https://github.com/aien-dev/omega/tree/614393d5b3f3d9d06f214c5f9a30a5788166bb07/src/runtime}
\end{quote}

The central pattern is

\[
\text{publication}
\rightarrow
\text{dependency becomes ready}
\rightarrow
\text{reaction executes}
\rightarrow
\text{verified publication}
\rightarrow
\text{next dependency}.
\]

The resulting causal system was qualified in the R13 living-system experiment.

The actual machine receipt is public:

\begin{quote}
\url{https://github.com/aien-dev/omega/blob/89e72ed72ba004631e1a0d81b874a136973d905f/evidence/R13/48d5a36c0ec2a8885aa21ebe1fd9d5d2259db1a25e935749adde6a1addd3b5b3.json}
\end{quote}

This was not a simulated architecture diagram. It ran on the GB10 system.

The human supplied a goal: reduce the cost of an Omega integer matrix-vector operation below \(55\%\) of the confirmed incumbent while preserving semantics.

The incumbent was

\[
7000\ \mathrm{ns/call}.
\]

The target was

\[
3850\ \mathrm{ns/call}.
\]

The realized candidate selected by the system measured

\[
2{,}686{,}162\ \mathrm{ps/call}
\approx
2686.162\ \mathrm{ns/call}.
\]

The receipt records:

\begin{itemize}
    \item \(137\) AIEN activations;
    \item \(122\) Omega activations;
    \item \(16\) GPU claims;
    \item \(16\) GPU completions;
    \item generation \(1\rightarrow2\);
    \item selected realization identity exactly matching the promoted identity;
    \item the same identity reconstructed from durable bytes;
    \item \(3{,}159\) causal records;
    \item \(11\) durable provenance links verified;
    \item zero production results using the unpromoted realization;
    \item all five control conditions passing;
    \item final AIEN goal status \texttt{MET}.
\end{itemize}

I do not call this proof of general intelligence.

It proves something narrower and more useful.

I built a system in which a goal caused a chain of planning, search, verification, physical experimentation, evidence publication, selection, promotion, and reassessment without a central program manually sequencing those stages.

That is an implementation of the process I describe as

\[
\text{conjecture}
\rightarrow
\text{test}
\rightarrow
\text{retain}
\rightarrow
\text{continue}.
\]

The architecture therefore existed before I had a general numerical measure for its epistemic progress.

\section{A System That Cannot Survive Being Wrong Is Not Continuing}

The next test was more important to me than another successful optimization.

I deliberately attacked the system.

The R14 qualification injected failures into the same resident architecture.

The receipt is public:

\begin{quote}
\url{https://github.com/aien-dev/omega/blob/2c647c6ff54be87c96041ce1fdd0abbdd52a19a9/evidence/R14/0c09168743bfc8acac8b137a561958e9644c1d4160f92400b9a4bcec7b18965f.json}
\end{quote}

Six classes of failure were tested.

\subsection{Corrupt candidate}

Candidate realizations were tampered with, made semantically incorrect, and made to crash.

The verifier refused them.

The bad realization was never promoted.

The valid candidate remained available.

\subsection{Forged authority}

The test presented forged capability references, valid references under the wrong subject, revoked references, stale generations, forged GPU claims, stuffed grant slots, forged verification records attributed to AEGIS, and a promotion forgery.

The receipt records them blocked or contained.

There were

\[
0
\]

publications under the stuffed grant and

\[
0
\]

mints during the attacks.

\subsection{Reaction cycles}

I injected

\[
A\leftrightarrow B
\]

and

\[
X\rightarrow Y\rightarrow Z\rightarrow X,
\]

plus a self-oscillating reaction.

They were bounded by the per-episode activation limit and oscillation containment.

Maximum cycle activations were

\[
256.
\]

Containment occurred in approximately

\[
6\ \mathrm{ms}.
\]

The legitimate organism continued producing work while the pathological reactions were contained.

\subsection{GPU saturation}

Sixteen lanes each submitted \(200\) requests against a four-claim resource budget.

The receipt records:

\begin{itemize}
    \item \(3200\) submitted;
    \item \(3200\) completed;
    \item \(3200\) checked;
    \item zero duplicates;
    \item zero lost;
    \item zero mismatches.
\end{itemize}

\subsection{GPU worker destruction}

The running GPU channel was deliberately destroyed while it held a claim.

The claim failed once.

Its resource charge was released.

The seat generation advanced.

The stale result was refused.

The retry was causally linked to the loss and published once.

\subsection{Crash during promotion}

The system was crashed at eight different points of the durable-generation promotion protocol.

Recovery returned either the complete old generation or the complete new generation.

It never produced a torn generation.

The full R14 run produced

\[
180{,}808
\]

causal records with causal verification passing.

These are causal reaction records, not crumbs in the Crumb Doctrine sense: a Doctrine crumb is a verified program plus its verification receipt.

Just as importantly, the experiment found real defects in earlier stages of the architecture. Those defects were fixed at the layer that owned them and preserved in regression tests.

That is not an embarrassment in this design.

It is the point.

An architecture claiming to learn from evidence but rewriting away its failures would contradict the ideal it is supposed to implement.

\section{I Required the Architecture to Beat a Control}

There is another failure mode in systems research: building an elegant architecture and assuming it must therefore be better.

I did not want that.

The R15 performance specification states:

\begin{quote}
\emph{Do not accept architecture on aesthetic grounds.}
\end{quote}

The preregistered methodology is public:

\begin{quote}
\url{https://github.com/aien-dev/omega/blob/4d47ce7fe85c3dc6bcfa5f5157250194656b3154/spec/r15-performance-proof.md}
\end{quote}

I built a sequential control using the same semantic work, same objects, same authority, same evidence, and same GPU seat.

Only the orchestration model changed.

That gave me a direct comparison between

\[
\text{resident causal architecture}
\]

and

\[
\text{sequential central orchestration}.
\]

The final qualification is public:

\begin{itemize}
    \item \url{https://github.com/aien-dev/omega/blob/de4f32dee665b45faa7f96cb87015830e2b7f313/evidence/R15/ATTEMPT-2-PASS.md}
    \item \url{https://github.com/aien-dev/omega/blob/de4f32dee665b45faa7f96cb87015830e2b7f313/evidence/R15/065c688408aa18950f363b52be2357c3cb853827592dc71ea72b6da181b4569a.json}
\end{itemize}

The qualification ran on the DGX Spark GB10 and passed

\[
16/16
\]

gates.

One measured difference was CPU to GPU synchronization:

\[
3.0
\]

synchronization events per GPU result for the one-worker resident configuration versus

\[
4.0
\]

for the sequential control.

The methodology also measures semantic operations per second, activation latency, dependency propagation, scheduler overhead, verification fast-path cost, generation barriers, serialization, CPU to GPU synchronization, memory traffic, GPU residency, wasted reactions, conflict and invalidation, causal-evidence overhead, and energy per semantic result.

That last metric is directly relevant to the Turing.

I had already reached the point where I was asking:

\begin{quote}
What did this semantic result physically cost?
\end{quote}

The missing half was:

\begin{quote}
How much epistemic gain did the semantic result represent?
\end{quote}

\section{The Architecture Also Needed to Admit When My Own Idea Was Worse}

The strongest evidence for the philosophy of the project may be a failed experiment.

I built the first TURING Field as an evidence-backed decision-record layer around multiple realizations of one semantic operation.

The source and tests are public:

\begin{itemize}
    \item \url{https://github.com/aien-dev/omega/blob/cc341abbe203c64506b2acab824f2e58e43f99ab/src/turing/field.h}
    \item \url{https://github.com/aien-dev/omega/blob/cc341abbe203c64506b2acab824f2e58e43f99ab/tests/turing/test_turing.c}
    \item \url{https://github.com/aien-dev/omega/blob/cc341abbe203c64506b2acab824f2e58e43f99ab/docs/turing/TURING_W0_PROPOSAL.md}
\end{itemize}

I expected the new selection rule to perform at least as well as a simple history-based control.

It did not.

The preregistered comparison produced:

\begin{center}
\begin{tabular}{lrr}
\toprule
 & Mean regret & Max regret\\
\midrule
Field v0 & \(3.58\%\) & \(25.51\%\)\\
Control & \(0.01\%\) & \(0.07\%\)\\
\bottomrule
\end{tabular}
\end{center}

The quality gate failed.

The replayability gate passed.

I did not change the threshold.

I did not delete the result.

I did not tune the benchmark until my method won.

The failed selector remains in the repository so the failure can continue to be reproduced:

\begin{quote}
\url{https://github.com/aien-dev/omega/blob/cc341abbe203c64506b2acab824f2e58e43f99ab/src/turing/field_select_v0_retired.c}
\end{quote}

The live Field was reframed around what the evidence actually supported: immutable, replayable decision records that cite their measurements.

The current code says directly:

\begin{quote}
\emph{Data only, not a subsystem.}
\end{quote}

The current tests ingest \(720\) evidence rows, reconstruct realization identities, verify source receipts, detect tampering, and verify that decisions cite the evidence they used.

A one-byte alteration of a receipt causes verification to fail.

This failure changed the architecture.

The tie policy was retired and the architecture moved toward lowest measured cost as the reference selection rule.

That history is preserved in:

\begin{quote}
\url{https://github.com/aien-dev/aien-architecture}
\end{quote}

This is the implementation-level reason I take the phrase

\begin{quote}
\emph{reality decides}
\end{quote}

seriously.

It has already overruled me.

\section{The Measurement Problem Appeared From Inside the Architecture}

By this point, I had built several pieces of the same pattern.

Omega could discover reusable structure.

The reaction architecture could continue from a goal into experimentation and promotion.

The recovery architecture could survive rejected hypotheses, forged authority, saturation, crashes, and cycles.

The performance architecture could measure the physical cost of semantic work.

TURING could preserve the evidence supporting a decision and preserve a failed selection rule as negative evidence.

These systems all shared the same implicit idea:

\begin{quote}
Improvement means finding a structure that explains or accomplishes more while costing less, and proving that the improvement survives independent evidence.
\end{quote}

But there was no common unit.

That was the abstraction boundary where I found the missing quantity.

\section{The Turing}

Let \(P\) be a declared measurement profile.

Let:

\begin{itemize}
    \item \(B\) be a baseline explanatory model;
    \item \(M\) be a candidate explanatory model;
    \item \(D\) be sealed held-out observations.
\end{itemize}

For a computable predictive distribution,

\[
L(D\mid M)
=
-\sum_{t=1}^{n}
\log_2p_M(x_t\mid x_{<t}).
\]

This is the predictive log-loss expressed in bits, following the coding relationship developed by Shannon~\cite{shannon1948}.

I also charge the explanation for existing:

\[
L(M).
\]

Therefore,

\[
DL(M,D)
=
L(M)+L(D\mid M).
\]

I define the Turing gain as

\[
\boxed{
\mathcal{T}(M;B,D,P)
=
DL(B,D)-DL(M,D).
}
\]

The reporting unit is the \textbf{Turing}, \(T\):

\[
\boxed{
1T
=
1\text{ bit of net held-out explanatory advantage}
}
\]

under the declared measurement profile.

Raw totals may only be compared at fixed evidence volume. I therefore also report the \textbf{specific Turing gain}: \(\mathcal{T}\) divided by the number of held-out observations, in Turings per observation. A measurement report separates five quantities that must not be collapsed: the raw total \(T\), the specific gain, the candidate's quality, the baseline's strength, and the evidence volume.

A note on the name, because I checked: an unrelated proposal called the ``Turing Ratio''~\cite{masum2003} is a chess performance rating on a decibel scale. It shares a name fragment with this unit and nothing else. I also searched the recent literature for the closest neighbors to this proposal. The nearest is SuperARC~\cite{superarc2025}, which defines intelligence through explanatory compression but as a benchmark score, not a named unit, with no baseline net-gain and no energy accounting. The formal intelligence measures of Legg and Hutter~\cite{legghutter2007} and Chollet~\cite{chollet2019} propose no unit at all, and the recent MDL theory of machine understanding~\cite{zhangliu2025} likewise names none. To my knowledge, no prior work uses the term \textit{Turing} for the protocol-bound quantity defined here.

A score of

\[
+50{,}000T
\]

means that, after paying the encoding cost of the new explanation, it required \(50{,}000\) fewer bits than the baseline to account for observations it did not receive during construction or selection.

A negative value is valid:

\[
\mathcal{T}<0
\]

means the explanation did not pay for itself.

\section{Why This Is Not Merely Renaming a Bit}

A bit is a unit of information.

A Turing is the reporting unit of a constrained comparative measurement.

The measured quantity is

\[
\mathcal{T}(M;B,D,P).
\]

A result might be

\[
\mathcal{T}(M;B,D,P)
=
84{,}220T.
\]

That statement implies:

\begin{itemize}
    \item a declared baseline;
    \item a frozen candidate;
    \item a computable encoding convention;
    \item held-out observations;
    \item a sealed evaluation protocol;
    \item a comparison of total description length;
    \item a result that may be negative;
    \item evidence from which another person can recompute the outcome.
\end{itemize}

I do not claim to have discovered a new fundamental dimension of physics.

I am giving a name and protocol to a particular measured difference.

\section{What the Turing Does Not Measure}

A precise instrument for one quantity is stronger than an indefensible measure of everything. The thermometer is the analogy I keep: nobody asks a thermometer to measure pressure, and its authority comes from refusing to.

The Turing measures net held-out explanatory compression. It does not measure semantic grounding, causal direction, agency, intervention ability, social understanding, values, alignment, or consciousness. A high score means the candidate made unseen observations less surprising after paying its own cost. It does not mean the candidate understands in any broader sense, and I claim nothing of the kind.

This narrowing is also what separates the unit from benchmark culture. A benchmark ranks systems on fixed tasks; the score conflates the candidate, the baseline, and the evidence volume, and the leaderboard becomes the product. The Turing protocol reports those quantities separately, under a declared profile, with the baseline named and the held-out evidence sealed. Two measurements are comparable only under the same profile, and the report must say so.

\section{Why Memorization Does Not Get a Free Pass}

Suppose a system memorizes a giant table.

Its prediction error can approach zero on data represented by that table.

But the table itself has to be encoded.

The relevant quantity is not merely

\[
L(D\mid M).
\]

It is

\[
L(M)+L(D\mid M).
\]

If the table costs

\[
8{,}000{,}000\text{ bits}
\]

and saves

\[
3{,}000{,}000\text{ prediction bits},
\]

the result is strongly negative.

A small rule can therefore outperform a memorizer epistemically even when both fit the same observations.

This is the same structural principle already implemented by Omega Library Discovery, generalized from

\[
\text{instructions saved}
-
\text{abstraction cost}
\]

to

\[
\text{prediction bits saved on unseen evidence}
-
\text{explanation cost}.
\]

\section{Physics Does Not Disappear}

While thinking about explanatory compression, I initially considered whether a new unit might replace familiar electrical measurements.

That would have been a category error.

Volts, amperes, watts, and joules measure physical quantities.

Turings measure explanatory gain.

I now treat them as orthogonal axes.

FORGE already provides a substrate-neutral physical boundary.

The FORGE V2 specification is public:

\begin{quote}
\url{https://github.com/aien-dev/physics/blob/2da7b9d8d429abb590c4a12b1bae95eb6fd4461e/docs/FORGE_SUBSTRATE_V2_SPEC.md}
\end{quote}

Its stated purpose is that FORGE realizes semantic operations on physical devices.

The substrate description advertises evidence capabilities including:

\begin{itemize}
    \item duration;
    \item energy;
    \item temperature;
    \item measured error;
    \item repeats;
    \item hardware status;
    \item faults.
\end{itemize}

The implementation is here:

\begin{quote}
\url{https://github.com/aien-dev/physics/tree/5969159c383217178b7b524fff13869489c8d799/forge/v2}
\end{quote}

The informational quantity is

\[
T.
\]

The physical quantity remains

\[
J.
\]

Their ratio is

\[
\boxed{\frac{T}{J}},
\]

which I call \textbf{Turing yield}.

It is not energy converted into intelligence.

It is net explanatory gain measured per joule consumed by the declared computation.

\section{The Physical Instrumentation Already Exists}

This proposal is not based on the assumption that I might someday measure energy.

Omega's R15 harness already does.

The implementation includes:

\begin{quote}
\url{https://github.com/aien-dev/omega/blob/ac09e6247da619e5cc1f29a70fa4792f10b679f1/tests/runtime/r15_measure.c}
\end{quote}

The measurement layer uses physical telemetry including energy counters, GPU power/energy instrumentation, ARM performance counters, CPU cycles, retired instructions, bus accesses, cache misses, memory accesses, thread CPU time, and wall-clock time.

Omega's empirical cost model already includes an energy dimension:

\begin{quote}
\url{https://github.com/aien-dev/omega/blob/8290082331072f45fb0b3a60ca4f7654190fbfb6/src/runtime/rx_costmodel.h}
\end{quote}

A prediction contains an energy estimate and a request may specify an energy budget.

Thus the architecture had already learned to distinguish

\[
\text{semantic correctness}
\]

from

\[
\text{physical cost}.
\]

The Turing does not replace that machinery.

It completes it.

\section{Multiple GPU Processes Require Attribution, Not New Physics}

If workloads \(A\) and \(B\) run together, I cannot assume

\[
E(A+B)=E(A)+E(B).
\]

They may contend for compute units, memory bandwidth, caches, unified memory, power budget, thermal headroom, clocks, and synchronization resources.

I therefore define an interaction experiment.

Measure

\[
E_0,\quad E_A,\quad E_B,\quad E_{AB}.
\]

Then define

\[
\boxed{
I_{AB}
=
E_{AB}-E_A-E_B+E_0.
}
\]

That interaction term must remain visible.

I do not allow an accounting system to pretend every shared joule belongs exclusively to one process.

The physical conservation requirement is

\[
E_{\text{measured}}
\approx
E_{\text{exclusive}}
+
E_{\text{shared}}
+
E_{\text{interaction}}
+
E_{\text{unattributed}}.
\]

If attributed energy exceeds what the machine measured outside the declared uncertainty, the attribution fails.

Likewise, I impose informational conservation.

If three processes jointly produce one explanation worth

\[
10{,}000T,
\]

I cannot award all three processes \(10{,}000T\) and announce \(30{,}000T\).

The explanatory gain itself cannot be duplicated.

\section{Observation Is Not Authority}

Measurement must not become control.

ARGUS is the defensive observation plane in AIENOS.

Its ABI is public:

\begin{quote}
\url{https://github.com/aien-dev/aienos/blob/b908e0187c2af7ab5675bafaf5768dc2f877441d/native/argus/argus_abi.h}
\end{quote}

Its boundary is explicit: ARGUS observes, detects, contains by request, and defends; it never authorizes.

A finding is evidence.

A containment request is a proposal.

Authority remains elsewhere.

That principle also constrains the Turing.

A system does not receive permission because it has a high \(T\).

It does not receive permission because an experiment has high expected information gain.

It does not get to promote itself merely because it claims to have learned something.

The measurement layer measures.

Authority remains authority.

\section{The First Turing Measurement Is Recorded}

I want to be precise about what I have and have not demonstrated.

On 2026-09-29 the first profile-bound idealized Turing result was recorded under measurement profile V0 (\href{https://github.com/aien-dev/omega/pull/85}{omega pull request~85}, merged the same day). The candidate was a table conditioning on the tried operation plus the last four outcomes of Omega search traces; the baseline was a table conditioning only on the last outcome. On sealed held-out seeds 8, 9, and 10, about 5.8~million search events never seen during model selection, the candidate reduced total held-out description length by \textbf{2,559,679.825 bits}, scoring \textbf{+2,559,679.825~$T_{\mathrm{ideal}}$} under profile V0, about \textbf{0.44~$T_{\mathrm{ideal}}$ per held-out observation} of specific Turing gain. The preregistered bar of 234,602~bits was cleared, with each held-out seed passing on its own. The scoring tool refuses to score the held-out seeds a second time, and a refit check refits both models from the training files and reproduces the receipt exactly. The complexity charge penalizes stored information that does not earn its cost on held-out observations. Under Lemma~2's assumptions, pure memorization cannot outperform the cheaper baseline. All scoring arithmetic uses fixed-point integer micro-bits, with worst-case rounding under 6~bits over the whole held-out set.

I state the result's status plainly, because the ledger now makes it checkable. TY-1+TY-2 is an internally replayed descriptive result under historical profile V0. The experiment record reports separation of fitting, selection, and evaluation seeds, but independent blinding was not established. The original corpus is not publicly downloadable, so public end-to-end replay is not currently available. Statistical uncertainty, baseline sensitivity, and attribution ablations remain pending. The result therefore does not constitute full qualification under the current protocol. This does not establish contamination and does not make the arithmetic worthless; it distinguishes an internally reported held-out result from an independently auditable demonstration. The full evidence ledger is in the formal paper.

The honest limits, stated with the result: \(\mathcal{T}\) here is a theoretical code length. No entropy coder was built that produces the smaller file, so this is not yet a physically realized encoding. Energy was not measured in this run, so no Turing yield (\(T/J\)) is reported. Joining R15 energy measurements to Turing-gain receipts is the next step.

I will not score old runs retroactively.

The R13 result predates the Turing protocol.

The M11 abstraction experiment uses instruction-count compression, not probabilistic held-out description length.

Likewise, retirement of all remaining central orchestration is still an active architectural gate.

The architecture is therefore evidence for the \emph{implementation of the ideal}, and the TY-1+TY-2 result is a completed pre-EXP instrument result under profile V0. EXP-001's preregistered runs, which close the reference-coder bridge (H1: \(T_{\mathrm{ideal}} \approx T_{\mathrm{real}}\) within the declared bound), are pending; full qualification of the metric, including the physical coder bridge and energy-joined receipts, is still ahead.

\section{The First Scored Experiment}

The first scored experiment was intentionally boring.

I did not need a giant language model.

I already had an appropriate experimental substrate inside Omega: real search and synthesis traces, divided before candidate construction into training, validation, and sealed test.

The scored run compared two predictors: the candidate explanatory table (conditioning on the tried operation plus the last four outcomes) against the declared baseline (a table conditioning only on the last outcome). The wider comparison set, uniform, empirical frequency, a small Markov-style table, the existing fixed heuristic, belongs to the EXP-001 design, whose preregistered runs are pending. For every candidate \(M\), the protocol calculated \(L(M)\) and \(L(D\mid M)\), then \(\mathcal{T}(M;B,D,G,P)\) against the declared baseline.

The first gate was simple:

\[
\boxed{\mathcal{T}>0}
\]

It passed, as recorded in the previous section: +2,559,679.825~\(T_{\mathrm{ideal}}\) under profile V0, each held-out seed passing on its own. The 6-bit figure reported with the result is an arithmetic rounding bound on the scoring computation, not a statistical confidence interval; statistical uncertainty for \(\mathcal{T}\) is required future work.

One more discipline the protocol enforces, and I want to state it plainly because it is easy to get wrong: the cryptographic hash of the sealed test set proves \textit{which} data was held out. It does not prove the machine never saw it. A hash is a commitment to identity, not a blindfold. Non-access has to be established by a separate, auditable blinding mechanism, and the protocol treats the two claims separately: the hash proves the dataset is the one declared, the blinding proves the candidate was built without it. Either claim without the other is not a measurement.

In the formal paper the preregistered instrument-validation step is EXP-001. The scored run described above is not EXP-001; it is the completed pre-EXP instrument result TY-1+TY-2. EXP-001's preregistered runs are pending, and with them the reference-coder bridge: H1, that \(T_{\mathrm{ideal}} \approx T_{\mathrm{real}}\) within the declared bound, is open. EXP-001 is followed by EXP-002, Stochastic Law Discovery, in which a hidden evaluator generates quantized trajectories from an undisclosed stochastic law: pure diffusion, diffusion with drift, and mean-reverting dynamics, and the candidate must earn its complexity on sealed held-out trajectories without inventing deterministic structure inside noise; and by EXP-003, active experiment selection, in which the machine proposes the intervention that most efficiently distinguishes competing explanations. EXP-002 asks whether the instrument can detect discovered structure. That is the empirical ladder the measure is now climbing.

\section{The Second Experiment Joins the Machine Back to Physics}

After the informational measurement works independently, I will join it to the physical measurement.

The same work unit should produce:

\begin{itemize}
    \item candidate explanation;
    \item prediction trace;
    \item Turing receipt;
    \item physical interval receipt;
    \item energy attribution receipt.
\end{itemize}

Only then will I calculate

\[
T/s
\]

and

\[
T/J.
\]

For discovery itself I intend to distinguish evaluation cost from total discovery cost.

Define

\[
Y_{\mathrm{eval}}
=
\frac{T}{E_{\mathrm{eval}}}
\]

and

\[
\boxed{
Y_{\mathrm{discovery}}
=
\frac{T}
{
E_{\mathrm{search}}
+
E_{\mathrm{train}}
+
E_{\mathrm{synthesis}}
+
E_{\mathrm{verify}}
+
E_{\mathrm{eval}}
}.
}
\]

For the machine-scientist problem, the second quantity is more important. The discovery denominator must include the energy cost of training the candidate itself (\(E_{\mathrm{train}}\)): a model that discovers a compact explanation only after an enormous training expenditure has a lower discovery yield, and the measure should say so.

\section{Why This Matters}

I did not build this measure because I wanted a new number. I built it because three problems I care about have no honest number attached to them, and I think this one can serve all three without pretending to be something it is not.

First, the environment. As machine computation scales, the energy cost of producing understanding is invisible in every benchmark I know. Accuracy per watt tells you how cheaply a model performs; it does not tell you how cheaply a machine \textit{learns}. Turings per joule makes the energy cost of discovery legible. If the field ever has to answer how much understanding it bought per unit of energy, this is the shape of the answer.

Second, honesty. The protocol only works if the baseline is declared in advance, the held-out data is sealed, every artifact is recomputable, and negative results are reported as what they are: a candidate that did not pay for itself. That discipline is not a nicety. It is the difference between a measure and a marketing claim. A field that adopts this unit adopts, along with it, the habit of saying what it tried and what failed.

Third, alignment. A great deal of alignment work has no measurable target: we ask machines to be safe, helpful, and honest, and we grade them with vibes. Verified understanding, operationalized as net held-out explanatory compression, is one thing we can actually measure. It is not the whole of alignment, and I claim nothing more.

And the honest bounds, stated once so they cannot be mistaken for modesty: the Turing does not measure safety. It does not measure consciousness, intent, or moral worth. It measures whether a machine discovered an explanation that makes unseen reality less surprising, after paying the explanation's own cost. That is all it is, and that is enough to be worth measuring.

\section{Why I Think This Extends Turing's Method}

I am not proposing to replace the Turing Test.

The Turing Test asks an operational behavioral question: can a machine's responses be distinguished from those of a human under an interrogation protocol?~\cite{turing1950}

The question I encountered while building AIEN was different.

I already had machines capable of producing fluent responses.

That did not solve my problem.

I needed to know whether the system had discovered structure that generalized beyond what it had already seen.

So I arrived at a different operational question:

\begin{quote}
\textbf{Can the machine construct an explanation that pays for its own complexity by making unseen reality less surprising?}
\end{quote}

The progression I see is

\[
\text{imitation}
\rightarrow
\text{prediction}
\rightarrow
\text{explanation}
\rightarrow
\text{experiment}.
\]

Turing operationalized an argument about appearances.

I am trying to operationalize an argument about explanatory progress.

The intellectual debt is methodological, not rhetorical.

\section{The Architecture Is the Argument}

The reason I am willing to publish the idea now is not that the equation looks elegant.

It is that I can point to pieces of a working machine that forced the equation on me.

When I say \emph{compression}, I can point to Omega's implemented abstraction-discovery engine and its qualification receipt.

When I say \emph{generalization}, I can point to held-out search acceleration in the same experiment.

When I say \emph{continuation}, I can point to the R13 silicon receipt and its causal chain from human goal to new generation.

When I say \emph{falsification}, I can point to R14 fault injection and the defects those tests actually exposed.

When I say \emph{physical cost}, I can point to R15's measured energy, PMU, synchronization, and runtime instrumentation.

When I say \emph{negative evidence}, I can point to TURING Wave 1, where my proposed selector lost and remains permanently reproducible as a failed experiment.

When I say \emph{evidence before authority}, I can point to AEGIS certification, the capability system, ARGUS, and the promotion barrier.

When I say \emph{different physical realizations of the same meaning}, I can point to FORGE's substrate-neutral machine boundary and Omega's realization system.

This is why I do not describe the Turing as an idea I invented in isolation.

I discovered the need for it while working inside this architecture.

The abstraction itself exposed the missing measurement.

\section{Falsification}

I want this proposal to be able to die.

I therefore consider the following failures binding.

\subsection{F1: Prediction and compression diverge}

If better probabilistic prediction does not produce correspondingly smaller real lossless encodings under the declared coding procedure, the bridge fails for that regime.

\subsection{F2: No held-out gain}

If candidates compress training traces but do not produce positive net Turing gain on sealed held-out observations after paying \(L(M)\), I have demonstrated fitting rather than explanatory progress.

\subsection{F3: Profile instability}

If irrelevant implementation changes materially alter \(T\), the measurement profile is inadequate.

\subsection{F4: Physical accounting does not close}

If attributed energy does not reconstruct measured energy within uncertainty, I will not report \(T/J\).

Importantly,

\[
T
\]

may survive this failure, while

\[
T/J
\]

may not.

\subsection{F5: A better measure replaces it}

If another operational measure captures generalized explanatory progress more reproducibly and with fewer arbitrary conventions, I should use that instead.

The Turing should not receive protection from falsification merely because I named it.

\section{What I Am Actually Claiming}

My claim is narrower than ``I built intelligence.''

I am claiming that I built a computational architecture in which several properties normally discussed abstractly are already explicit engineering objects:

\begin{itemize}
    \item candidate explanations;
    \item semantic identity;
    \item multiple physical realizations;
    \item verification;
    \item evidence;
    \item authority;
    \item causal history;
    \item generation;
    \item failure;
    \item recovery;
    \item physical resource cost.
\end{itemize}

While working in that abstraction, I found that the missing epistemic quantity is naturally expressed as

\[
\text{cost of old explanation}
-
\text{cost of new explanation}.
\]

When evaluated on unseen observations and charged for the candidate itself, that becomes

\[
\boxed{
\mathcal{T}(M;B,D,P)
=
[L(B)+L(D\mid B)]
-
[L(M)+L(D\mid M)].
}
\]

I call the resulting unit the Turing.

The architecture does not yet prove that this is the right general measure of machine understanding.

It gives me something more useful:

\begin{quote}
\textbf{a machine on which the claim can actually be tested.}
\end{quote}

\section{Conclusion}

Alan Turing's enduring contribution to this problem was not merely a particular test.

It was a methodological move.

When a question becomes trapped in definitions, replace it with an operational question whose answer can be exposed to evidence.

I encountered the same problem from the other direction.

I built a machine around evidence.

I separated meaning from implementation.

I separated proposals from permission.

I required improvements to survive verification.

I made generations durable.

I made failures remain in the record.

I tested the system on real silicon.

I broke it deliberately.

I measured its physical costs.

And then I reached a question the architecture could not yet answer:

\begin{quote}
\textbf{How much explanatory ground did the machine actually gain?}
\end{quote}

That question produced the Turing.

Whether the Turing survives is now an experimental matter.

That is exactly how I want it.

\begin{center}
\textbf{The machine proposes.}\\
\textbf{The experiment measures.}\\
\textbf{Reality decides.}
\end{center}

\section{Epilogue}

After this engineering history was written, TY-1+TY-2 produced the first profile-bound idealized Turing result (+2,559,679.825 \(T_{\mathrm{ideal}}\)). EXP-001 remains the instrument-validation step required to close the reference-coder bridge (H1: \(T_{\mathrm{ideal}} \approx T_{\mathrm{real}}\) within the declared bound).

\section*{Reproducibility Manifest}

Each entry names the repository, the commit whose tree holds the cited artifact, the artifact SHA-256, the experiment it belongs to, the hardware profile where one was declared, and the receipt identifier where one exists. Where I could not establish a value from the repository or the receipt itself, I mark it [unavailable]; nothing is invented.


\begin{longtable}{p{0.22\textwidth}p{0.72\textwidth}}
\toprule
\textbf{Artifact} & \textbf{Reproducibility fields} \\
\midrule
\endfirsthead
\toprule
\textbf{Artifact} & \textbf{Reproducibility fields} \\
\midrule
\endhead
\bottomrule
\endfoot
Omega repository &
Repository: \texttt{aien-dev/omega} \newline
Commit: {\footnotesize\texttt{2d52219708b42c5d9b2e8f7d577fbb5226631a78}} \newline
Artifact SHA-256: [unavailable] \newline
Experiment ID: [unavailable] \newline
Hardware profile: [unavailable] \newline
Receipt ID: [unavailable] \newline
Public location: \url{https://github.com/aien-dev/omega} \\
\midrule
Omega Library Discovery source &
Repository: \texttt{aien-dev/omega} \newline
Commit: {\footnotesize\texttt{4f55f45ed13f63cf8171bebf223c6258f8351ab9}} \newline
Artifact SHA-256: {\footnotesize\texttt{2adebb01bb69c2c9\allowbreak{}36fdae73ef97dcdd\allowbreak{}45aaa965a854713a\allowbreak{}05ab831b42f701e4}} \newline
Experiment ID: M11 \newline
Hardware profile: [unavailable] \newline
Receipt ID: [unavailable] \newline
Public location: \url{https://github.com/aien-dev/omega/blob/4f55f45ed13f63cf8171bebf223c6258f8351ab9/src/omega_discovery.c} \\
\midrule
Omega Library Discovery header &
Repository: \texttt{aien-dev/omega} \newline
Commit: {\footnotesize\texttt{1ee43fbf1c7d588734b80bd0743e5d054aea750b}} \newline
Artifact SHA-256: {\footnotesize\texttt{4d607ec8f6ecb45c\allowbreak{}3f12cc1b0923fd09\allowbreak{}a8dc92bb9b7983c8\allowbreak{}c200d7f2555338b0}} \newline
Experiment ID: M11 \newline
Hardware profile: [unavailable] \newline
Receipt ID: [unavailable] \newline
Public location: \url{https://github.com/aien-dev/omega/blob/1ee43fbf1c7d588734b80bd0743e5d054aea750b/src/omega_discovery.h} \\
\midrule
M11 qualification receipt &
Repository: \texttt{aien-dev/omega} \newline
Commit: {\footnotesize\texttt{8e0a50a45814c5e3b427efd36a83a1e2e79a7500}} \newline
Artifact SHA-256: {\footnotesize\texttt{193bfce34119a66d\allowbreak{}7b8b739347a9ddd9\allowbreak{}10b91348359ef59f\allowbreak{}61d3f2cad579638d}} \newline
Experiment ID: M11 \newline
Hardware profile: Linux, aarch64 (per receipt host record) \newline
Receipt ID: CONTRACT-OMEGA-LIBRARY-DISCOVERY-M11 \newline
Public location: \url{https://github.com/aien-dev/omega/blob/8e0a50a45814c5e3b427efd36a83a1e2e79a7500/evidence/omega_library_discovery_qualification_receipt.json} \\
\midrule
R13 living-system silicon receipt &
Repository: \texttt{aien-dev/omega} \newline
Commit: {\footnotesize\texttt{89e72ed72ba004631e1a0d81b874a136973d905f}} \newline
Artifact SHA-256: {\footnotesize\texttt{48d5a36c0ec2a888\allowbreak{}5aa21ebe1fd9d5d2\allowbreak{}259db1a25e935749\allowbreak{}adde6a1addd3b5b3}} \newline
Experiment ID: R13 \newline
Hardware profile: DGX Spark (NVIDIA GB10), node spark-b87b, Linux 7.0.0-1019-nvidia, aarch64; silicon run \newline
Receipt ID: 20260928T122742Z-83f880786202 \newline
Run code commit: {\footnotesize\texttt{83f880786202bb2274848933395d55295cb0b679}} \newline
Public location: \url{https://github.com/aien-dev/omega/blob/89e72ed72ba004631e1a0d81b874a136973d905f/evidence/R13/48d5a36c0ec2a8885aa21ebe1fd9d5d2259db1a25e935749adde6a1addd3b5b3.json} \\
\midrule
R14 recovery qualification receipt &
Repository: \texttt{aien-dev/omega} \newline
Commit: {\footnotesize\texttt{2c647c6ff54be87c96041ce1fdd0abbdd52a19a9}} \newline
Artifact SHA-256: {\footnotesize\texttt{0c09168743bfc8ac\allowbreak{}ac8b137a561958e9\allowbreak{}644c1d4160f92400\allowbreak{}b9a4bcec7b18965f}} \newline
Experiment ID: R14 \newline
Hardware profile: DGX Spark (NVIDIA GB10), node spark-b87b, Linux 7.0.0-1019-nvidia, aarch64; silicon run \newline
Receipt ID: 20260928T134432Z-6b38173cd5be \newline
Run code commit: {\footnotesize\texttt{6b38173cd5bee75cbdb0b132d535ec0f0df38d4c}} \newline
Public location: \url{https://github.com/aien-dev/omega/blob/2c647c6ff54be87c96041ce1fdd0abbdd52a19a9/evidence/R14/0c09168743bfc8acac8b137a561958e9644c1d4160f92400b9a4bcec7b18965f.json} \\
\midrule
R15 preregistered methodology &
Repository: \texttt{aien-dev/omega} \newline
Commit: {\footnotesize\texttt{4d47ce7fe85c3dc6bcfa5f5157250194656b3154}} \newline
Artifact SHA-256: {\footnotesize\texttt{cae61c9bb0a1e48f\allowbreak{}b316e8159a296f50\allowbreak{}3ce155015eec3a02\allowbreak{}c814478f62dfda0a}} \newline
Experiment ID: R15 \newline
Hardware profile: [unavailable] \newline
Receipt ID: [unavailable] \newline
Public location: \url{https://github.com/aien-dev/omega/blob/4d47ce7fe85c3dc6bcfa5f5157250194656b3154/spec/r15-performance-proof.md} \\
\midrule
R15 qualification summary &
Repository: \texttt{aien-dev/omega} \newline
Commit: {\footnotesize\texttt{de4f32dee665b45faa7f96cb87015830e2b7f313}} \newline
Artifact SHA-256: {\footnotesize\texttt{5b3541e9d3f7568e\allowbreak{}5e74c5429e9dc0e5\allowbreak{}ecec297c5626b9f6\allowbreak{}6d01749bbd401172}} \newline
Experiment ID: R15 \newline
Hardware profile: [unavailable] \newline
Receipt ID: [unavailable] \newline
Public location: \url{https://github.com/aien-dev/omega/blob/de4f32dee665b45faa7f96cb87015830e2b7f313/evidence/R15/ATTEMPT-2-PASS.md} \\
\midrule
R15 machine-readable receipt &
Repository: \texttt{aien-dev/omega} \newline
Commit: {\footnotesize\texttt{de4f32dee665b45faa7f96cb87015830e2b7f313}} \newline
Artifact SHA-256: {\footnotesize\texttt{065c688408aa1895\allowbreak{}0f363b52be2357c3\allowbreak{}cb853827592dc71e\allowbreak{}a72b6da181b4569a}} \newline
Experiment ID: R15 \newline
Hardware profile: DGX Spark GB10, host spark-b87b, kernel 7.0.0-1019-nvidia; silicon mode \newline
Receipt ID: 20260929T025735Z-3e9e53be3358-silicon \newline
Run code commit: {\footnotesize\texttt{3e9e53be3358f6f2849a4f02044a1fd809240cbb}} \newline
Public location: \url{https://github.com/aien-dev/omega/blob/de4f32dee665b45faa7f96cb87015830e2b7f313/evidence/R15/065c688408aa18950f363b52be2357c3cb853827592dc71ea72b6da181b4569a.json} \\
\midrule
TURING Field record definitions &
Repository: \texttt{aien-dev/omega} \newline
Commit: {\footnotesize\texttt{cc341abbe203c64506b2acab824f2e58e43f99ab}} \newline
Artifact SHA-256: {\footnotesize\texttt{62f5eae7c870ce87\allowbreak{}c5b3a20b0d99473c\allowbreak{}c2bc0b5eb617398c\allowbreak{}0555ef08ad22823c}} \newline
Experiment ID: TURING W0 \newline
Hardware profile: [unavailable] \newline
Receipt ID: [unavailable] \newline
Public location: \url{https://github.com/aien-dev/omega/blob/cc341abbe203c64506b2acab824f2e58e43f99ab/src/turing/field.h} \\
\midrule
TURING preregistered experiment &
Repository: \texttt{aien-dev/omega} \newline
Commit: {\footnotesize\texttt{cc341abbe203c64506b2acab824f2e58e43f99ab}} \newline
Artifact SHA-256: {\footnotesize\texttt{a709f74fdadcf5f7\allowbreak{}20a5bfb62c8949e6\allowbreak{}cd439363c676cf25\allowbreak{}526147ef65b08705}} \newline
Experiment ID: TURING W0 \newline
Hardware profile: [unavailable] \newline
Receipt ID: [unavailable] \newline
Public location: \url{https://github.com/aien-dev/omega/blob/cc341abbe203c64506b2acab824f2e58e43f99ab/docs/turing/TURING_W0_PROPOSAL.md} \\
\midrule
TURING reproducibility and tamper tests &
Repository: \texttt{aien-dev/omega} \newline
Commit: {\footnotesize\texttt{cc341abbe203c64506b2acab824f2e58e43f99ab}} \newline
Artifact SHA-256: {\footnotesize\texttt{50c9a8174f19b348\allowbreak{}d101d3dd352d1c99\allowbreak{}25b5fa1a688a44e0\allowbreak{}8572aabfd10af54e}} \newline
Experiment ID: TURING W0 \newline
Hardware profile: [unavailable] \newline
Receipt ID: [unavailable] \newline
Public location: \url{https://github.com/aien-dev/omega/blob/cc341abbe203c64506b2acab824f2e58e43f99ab/tests/turing/test_turing.c} \\
\midrule
Omega empirical cost model &
Repository: \texttt{aien-dev/omega} \newline
Commit: {\footnotesize\texttt{8290082331072f45fb0b3a60ca4f7654190fbfb6}} \newline
Artifact SHA-256: {\footnotesize\texttt{89b9b953099e10eb\allowbreak{}9a3ce01ec109d72c\allowbreak{}18153f67f3f41399\allowbreak{}35c8997e2a6cffc1}} \newline
Experiment ID: R15 \newline
Hardware profile: [unavailable] \newline
Receipt ID: [unavailable] \newline
Public location: \url{https://github.com/aien-dev/omega/blob/8290082331072f45fb0b3a60ca4f7654190fbfb6/src/runtime/rx_costmodel.h} \\
\midrule
FORGE V2 specification &
Repository: \texttt{aien-dev/physics} \newline
Commit: {\footnotesize\texttt{2da7b9d8d429abb590c4a12b1bae95eb6fd4461e}} \newline
Artifact SHA-256: {\footnotesize\texttt{a9c15f8b2dc52810\allowbreak{}6b1690e6147fb578\allowbreak{}3c090648e828397c\allowbreak{}95c007699f6458d3}} \newline
Experiment ID: [unavailable] \newline
Hardware profile: [unavailable] \newline
Receipt ID: [unavailable] \newline
Public location: \url{https://github.com/aien-dev/physics/blob/2da7b9d8d429abb590c4a12b1bae95eb6fd4461e/docs/FORGE_SUBSTRATE_V2_SPEC.md} \\
\midrule
FORGE V2 implementation &
Repository: \texttt{aien-dev/physics} \newline
Commit: {\footnotesize\texttt{5969159c383217178b7b524fff13869489c8d799}} \newline
Artifact SHA-256: [unavailable] \newline
Experiment ID: [unavailable] \newline
Hardware profile: [unavailable] \newline
Receipt ID: [unavailable] \newline
Public location: \url{https://github.com/aien-dev/physics/tree/5969159c383217178b7b524fff13869489c8d799/forge/v2} \\
\midrule
ARGUS ABI &
Repository: \texttt{aien-dev/aienos} \newline
Commit: {\footnotesize\texttt{b908e0187c2af7ab5675bafaf5768dc2f877441d}} \newline
Artifact SHA-256: {\footnotesize\texttt{7368f92744b0d048\allowbreak{}1b64462fa00f873f\allowbreak{}0820e7a612171d9c\allowbreak{}65d476222e952e51}} \newline
Experiment ID: [unavailable] \newline
Hardware profile: [unavailable] \newline
Receipt ID: [unavailable] \newline
Public location: \url{https://github.com/aien-dev/aienos/blob/b908e0187c2af7ab5675bafaf5768dc2f877441d/native/argus/argus_abi.h} \\
\midrule
AIEN architecture &
Repository: \texttt{aien-dev/aien-architecture} \newline
Commit: {\footnotesize\texttt{33acd896e5d0cebe18201472b8f6b0b2e91f1b91}} \newline
Artifact SHA-256: [unavailable] \newline
Experiment ID: [unavailable] \newline
Hardware profile: [unavailable] \newline
Receipt ID: [unavailable] \newline
Public location: \url{https://github.com/aien-dev/aien-architecture} \\
\midrule
\end{longtable}


\appendix
\section{Revision history}

This companion paper was revised through v1.7 on 2026-10-02 before release. The revision record is kept here, outside the body of the paper, for transparency.

\begin{description}
\item[v1.7] Result-status reconciliation with the formal preprint v1.13: result-status paragraph added beside the TY-1+TY-2 measurement (internally replayed descriptive result; independent blinding not established; corpus not publicly downloadable); ``sealed unseen'' replaced by ``sealed held-out'' with the blinding distinction; cross-reference updated to the formal preprint v1.13.
\item[v1.5] Terminology and status reconciliation with the formal preprint v1.9: ``first canonical Turing measurement'' replaced throughout by the first profile-bound idealized Turing result, +2,559,679.825~\(T_{\mathrm{ideal}}\) under profile V0; the TY-1+TY-2 run restated as the completed pre-EXP instrument result, with EXP-001's preregistered runs pending and H1 open; cross-reference updated to the formal preprint v1.9.
\item[v1.1] Cross-reference updated to the formal preprint v1.3; forward pointer to the preregistered experiment ladder EXP-001/002/003 defined there.
\item[v1.2] Em dashes removed throughout at the author's direction (punctuation only, no content change).
\item[v1.3] Cross-reference updated to the formal preprint v1.6; first canonical Turing measurement recorded (omega PR 85, profile V0); AEGIS role clarified per the AEGIS ruling; FORGE terminology; narrow-scope statement and specific Turing gain.
\item[v1.4] Epilogue added (TY-1+TY-2 idealized Turing result; EXP-001 as the reference-coder bridge); cross-reference updated to the formal preprint v1.8; live branch URLs replaced with immutable commit permalinks; evidence index rebuilt as a reproducibility manifest (repository, commit, artifact SHA-256, experiment ID, hardware profile, receipt ID); R13/R14 causal records renamed out of the Crumb Doctrine vocabulary with the distinction stated; canonical role vocabulary applied throughout (AIEN proposes, Omega defines semantics and synthesizes candidates, AEGIS authorizes/certifies, independent verifiers verify, PHYSICS/FORGE realizes and measures, hardware executes, evidence records, Cortex learns/retains).
\end{description}

\begin{thebibliography}{9}

\bibitem{turing1950}
A.~M. Turing,
\newblock ``Computing Machinery and Intelligence,''
\newblock \emph{Mind}, vol.~59, no.~236, pp.~433--460, 1950.

\bibitem{shannon1948}
C.~E. Shannon,
\newblock ``A Mathematical Theory of Communication,''
\newblock \emph{Bell System Technical Journal}, vol.~27,
pp.~379--423 and 623--656, 1948.

\bibitem{rissanen1978}
J. Rissanen,
\newblock ``Modeling by Shortest Data Description,''
\newblock \emph{Automatica}, vol.~14, no.~5, pp.~465--471, 1978.

\bibitem{grunwald2007}
P. Gr{\"u}nwald,
\newblock \emph{The Minimum Description Length Principle},
\newblock MIT Press, 2007.

\bibitem{solomonoff1964}
R.~J. Solomonoff,
\newblock ``A Formal Theory of Inductive Inference,''
\newblock \emph{Information and Control}, vol.~7, parts I--II, 1964.

\bibitem{kolmogorov1965}
A.~N. Kolmogorov,
\newblock ``Three Approaches to the Quantitative Definition of Information,''
\newblock \emph{Problems of Information Transmission}, vol.~1, no.~1, pp.~1--7, 1965.

\bibitem{haldane1949}
J.~B.~S. Haldane,
\newblock ``Suggestions as to Quantitative Measurement of Rates of Evolution,''
\newblock \emph{Evolution}, vol.~3, pp.~51--56, 1949.

\bibitem{masum2003}
H. Masum, S. Christensen, and F. Oppacher,
\newblock ``The Turing Ratio: A Framework for Open-Ended Task Metrics,''
\newblock \emph{Journal of Evolution and Technology}, vol.~13, no.~2, 2003.

\bibitem{superarc2025}
A. Hern{\'a}ndez-Espinosa, F. Ozelim, F. Abrah{\~a}o, and H. Zenil,
\newblock ``SuperARC: A Test of Machine Superintelligence,''
\newblock arXiv:2503.16743, 2025.

\bibitem{legghutter2007}
S. Legg and M. Hutter,
\newblock ``Universal Intelligence: A Definition of Machine Intelligence,''
\newblock \emph{Minds and Machines}, vol.~17, no.~4, pp.~391--444, 2007.

\bibitem{chollet2019}
F. Chollet,
\newblock ``On the Measure of Intelligence,''
\newblock arXiv:1911.01547, 2019.

\bibitem{zhangliu2025}
Canlin Zhang and Xiuwen Liu,
\newblock ``A Theory of Machine Understanding via the Minimum Description Length Principle,''
\newblock arXiv:2504.00395, 2025.

\end{thebibliography}

\end{document}
